\uuid{ikTv}
\exo7id{2574}
\titre{exo7 2574}
\auteur{delaunay}
\organisation{exo7}
\datecreate{2009-05-19}
\isIndication{false}
\isCorrection{false}
\chapitre{Réduction d'endomorphisme, polynôme annulateur}
\sousChapitre{Diagonalisation}
\cpge{mpsi-mp2i.changements-bases.c, mpsi-mp2i.determinants.d, pcsi.calcul-matriciel-systemes-lineaires.d, ptsi.calcul-matriciel-systemes-lineaires.d, mp-mpi.reduction-endomorphismes.d, pc.reduction.c, psi.reduction-endomorphismes.c, pt.reduction-endomorphismes.c, mp-mpi.reduction-endomorphismes.e, pc.reduction.e, psi.reduction-endomorphismes.e, pt.reduction-endomorphismes.d, mp-mpi.equations-differentielles-lineaires.c}
\module{Algèbre}
\niveau{L2}
\difficulte{}

\contenu{
\texte{({\it 5 points})  Soit $A$ une matrice carr\'ee, $A\in M_n(K)$ ($K=\R$ ou $\C$). On rappelle que la trace d'une matrice est la somme de ses coefficients diagonaux et que ${\rm tr}(BAB^{-1})={\rm tr}A$.

D\'emontrer que $\det(\exp A)={\rm e^{trA}}$ dans les cas suivants : 
}
\begin{enumerate}
    \item \question{$A$ diagonalisable.}
    \item \question{$A$ triangulaire sup\'erieure ayant une diagonale de z\'eros.}
    \item \question{$A$ trigonalisable.}
    \item \question{$A$ quelconque.
}
\end{enumerate}
}
